(1) 单位渠宽体积流率 $q=h_1u_1=h_2u_2$ 守恒。积分定常动量方程时，单位宽静水压力合力为 $\rho gh^2/2$，故 $q^2/h_1+gh_1^2/2=q^2/h_2+gh_2^2/2$。对 $h_2\ne h_1$ 整理得 $2q^2/g=h_1h_2(h_1+h_2)$。

(2) 令 $r=h_2/h_1$，用 $q^2/(gh_1^3)=Fr_1^2$ 得 $r(r+1)=2Fr_1^2$，所以 $\boxed{r=(\sqrt{1+8Fr_1^2}-1)/2}$。当 $Fr_1>1$，$r>1$，且 $u_2=u_1/r<u_1$。另一负根不代表正水深。

(3) 比能水头 $H=h+q^2/(2gh^2)$。两侧相减，$H_1-H_2=(h_1-h_2)+\frac{q^2}{2g}(h_1^{-2}-h_2^{-2})$。代入 $q^2/g=h_1h_2(h_1+h_2)/2$ 后化简为 $\boxed{H_1-H_2=(h_2-h_1)^3/(4h_1h_2)}>0$。因此水跃满足质量和动量守恒，但机械能转为湍动、热和小尺度漩涡。

(4) $Fr_2^2=u_2^2/(gh_2)=Fr_1^2/r^3=(r+1)/(2r^2)<1$ 对 $r>1$ 成立。$Fr_1\downarrow1$ 时 $r\downarrow1$，故跃前后深度差趋零；水头损失按 $(h_2-h_1)^3$ 消失，体现弱水跃的立方耗散。
